% 4.6 =============== Decentralized Equilibrium ===============

\section{均衡概念的运用：分散经济下的最优增长模型}  \label{sec-decentralized-equilibrium}

回顾本章开头的问题：我们希望证明，在一个分散经济最优增长模型中，其竞争均衡是帕累托最优的。



% 4.6.1 ===== AD均衡-NGM
\subsection{AD}

为了更清晰的阐明宏观经济学中相关的均衡概念，现在先考虑如下问题：
当模型中包含企业（firm）和消费者（consumer）在市场上的互动时，经济中的竞争均衡（Competitive Equilibrium）是怎样的？
假设消费者拥有所有的生产要素（资本）并将其租赁给企业；假设代表性家庭拥有这些企业，索取其利润；假设产品市场和要素市场（劳动与资本）
是完全竞争的（perfectly competitive），这意味着家庭和企业都将价格视为给定（take as given）。


% 4.6.2 ===== SQ均衡-NGM
\subsection{SQ}


% 4.6.3 ===== RCE-NGM
\subsection{Recursive Comptetitive Equilibrium}

在一些宏观模型中，均衡可能并非帕累托最优，也可能难以通过序列形式或Arrow-Debreu形式求解，这时我们通常会定义递归竞争均衡（Recursive Comptetitive Equilibrium）。

不难发现，在对AD和ST进行介绍时，我们所解决问题的出发形式都是动态优化而不是此前介绍的递归形式（动态规划）。事实上，
在递归方法下同样也可以定义均衡概念。同时，RCE时ST的递归表示。


% 4.6.4 ==== 三种均衡概念的总结
\subsection{Summerize}